% Vietnamese translation for OI Public Library
% Source-PDF: IOI中国国家候选队论文/1999/来煜坤.pdf
\documentclass[11pt]{article}
\usepackage{oipl-vn}

\title{Nắm vững bản chất, vận dụng linh hoạt: bàn sâu về quy hoạch động}
\author{Lai Yukun}
\date{}

\begin{document}
\maketitle

\origtitle{Grasping the essence and applying flexibly: an in-depth discussion of dynamic programming}
\sourcepath{IOI China candidate team papers / 1999 / Lai Yukun.pdf}

\paragraph{Từ khóa.}
Quy hoạch động; xây dựng ý tưởng; cài đặt.

\section*{Tóm Tắt}

Bài viết này thảo luận nội dung cốt lõi và các đặc điểm cơ bản của tư tưởng
quy hoạch động; khảo sát phạm vi áp dụng của nó, cũng như cách suy nghĩ chung
khi xác lập không gian bài toán con và công thức truy hồi. Thông qua các ví
dụ, bài viết chỉ ra một số cách xử lý trong quá trình xác lập bài toán con:
tăng cường mệnh đề, điều chỉnh hợp lý các biến xác định trạng thái để giúp
lập phương trình quy hoạch động, và dùng tiền xử lý để quá trình quy hoạch
động dễ cài đặt hơn. Tiếp đó, bài viết phân tích vấn đề tràn bộ nhớ có thể
xuất hiện khi cài đặt quy hoạch động và một số cách khắc phục. Kết luận của
bài viết là: điểm then chốt của quy hoạch động vẫn nằm ở việc lập mô hình
toán học hiệu quả cho từng bài toán khác nhau, rồi vận dụng linh hoạt trên cơ
sở nắm chắc bản chất.

\section{Dẫn nhập}

Quy hoạch động là một tư tưởng thiết kế chương trình quan trọng và có giá trị
ứng dụng rộng rãi. Với khá nhiều bài toán, dùng tư tưởng quy hoạch động để
thiết kế thuật toán thường đạt hiệu quả thời gian cao; vì vậy, đối với những
bài có thể giải bằng quy hoạch động, đây thường là lựa chọn sáng suốt.

Những bài toán có thể giải bằng quy hoạch động thường là bài toán tối ưu, và
một phần cục bộ của lời giải tối ưu, hoặc của lời giải đặc thù, thường cũng là
lời giải tối ưu của bài toán con tương ứng dưới những điều kiện phù hợp. Đồng
thời, lời giải tối ưu của bài toán phải có liên hệ nhất định với lời giải tối
ưu của các bài toán con để có thể thiết lập quan hệ truy hồi. Nếu quan hệ này
khó thiết lập, tức là lời giải đặc thù của bài toán không chỉ phụ thuộc vào
lời giải đặc thù của bài toán con mà còn phụ thuộc vào các lời giải tổng quát
của bài toán con, thì một mặt ta khó ghi lại quá nhiều ``lời giải tổng quát'',
mặt khác hiệu suất truy hồi cũng rất thấp. Ngoài ra, để thể hiện hiệu quả
thời gian cao của quy hoạch động, các bài toán con nên chồng lặp với nhau:
nhiều bài toán khác nhau cùng dùng chung một số bài toán con. Nếu các bài toán
con không chồng lặp, thường nên dùng phương pháp khác, như chia để trị.

Nói chung, quy hoạch động có thể được cài đặt khá hiệu quả bằng hai cách. Cách
thứ nhất là ghi nhớ từ trên xuống: dùng đệ quy hoặc cách không đệ quy để quy
việc tìm lời giải tối ưu của bài toán về việc tìm lời giải tối ưu của các bài
toán con, rồi ghi lại các kết quả đã tính để chia sẻ. Cách thứ hai, cũng là
cách chủ yếu, là truy hồi từ dưới lên. Vì chi phí của cách này nhỏ hơn cách
thứ nhất nên nó được dùng phổ biến; các thảo luận dưới đây đều xét cách cài
đặt này. Quy hoạch động có hiệu quả cao vì trong khi liên tục giảm quy mô bài
toán, nó ghi lại lời giải một cách hữu hiệu, tránh hiện tượng giải đi giải lại
cùng một bài toán con. Do đó, nếu vận dụng đúng, hiệu suất thường cao hơn tìm
kiếm rất nhiều.

Tư tưởng quy hoạch động cho ta một hướng suy nghĩ khi giải các bài toán tối
ưu liên quan đến bài toán con chồng lặp: xét các bài toán con theo vòng lặp,
giảm quy mô bài toán rồi cuối cùng giải được bài toán ban đầu. Phạm vi bài
toán thích hợp với quy hoạch động rất rộng. Bản thân tư tưởng quy hoạch động
là quan trọng, nhưng điều quan trọng hơn là phân tích cụ thể từng bài toán:
xem bài toán có đủ điều kiện dùng quy hoạch động hay không, xác định không
gian bài toán con và công thức truy hồi, rồi cài đặt những thuật toán này trên
máy tính có bộ nhớ thông thường hữu hạn. Dưới đây, ta lần lượt bàn sâu hơn về
hai phương diện: xây dựng ý tưởng và cài đặt.

\section{Xây dựng ý tưởng giải bài bằng quy hoạch động}

Khi đối diện một bài toán và cân nhắc dùng quy hoạch động, điểm rất quan trọng
là phán đoán bài toán ấy có thể được giải hiệu quả bằng quy hoạch động hay
không. Khi xây dựng thuật toán quy hoạch động, ta thường phải xét các bài toán
con liên quan, tức không gian bài toán con, cách lập công thức truy hồi, và
cuối cùng là cài đặt thuật toán. Thực ra các quá trình này thường đan xen với
nhau. Không gian bài toán con và quan hệ truy hồi vốn liên hệ chặt chẽ: để lập
được quan hệ truy hồi hiệu quả, đôi khi cần điều chỉnh không gian bài toán
con; ngược lại, không gian bài toán con đã được xác định sơ bộ cũng có thể gợi
ý cho ta cách lập công thức truy hồi. Việc cuối cùng có thể dùng một công thức
truy hồi để nối bài toán với các bài toán con hay không cũng trở thành căn cứ
chính để phán đoán bài toán có thể dùng quy hoạch động hay không. Vì thế, tách
riêng từng phần để xét cô lập, hay cố chia quá trình suy nghĩ thành vài đoạn
cứng nhắc, vừa khó vừa không cần thiết. Hơn nữa, quy hoạch động không có một
mô hình toán học cố định; nó thay đổi theo từng bài toán, nên không thể quy
nạp thành một phương pháp ``vạn năng''. Dù vậy, với đa số bài toán, vẫn có một
hướng suy nghĩ cơ bản.

Trước hết, cần phân tích sơ bộ xem bài toán có khả năng giải bằng quy hoạch
động hay không. Nếu một bài toán khó xác định bài toán con, hoặc giữa lời giải
đặc thù của bài toán và lời giải đặc thù của bài toán con hoàn toàn không có
liên hệ, ta nên cân nhắc các phương pháp khác như tìm kiếm. Phán đoán sơ bộ là
cần thiết, vì nó giúp tránh lãng phí thời gian. Thông thường, một bài toán có
thể dùng quy hoạch động hiệu quả về cơ bản thỏa các đặc điểm sau:
\begin{enumerate}
  \item Lời giải tối ưu của bài toán con chỉ liên quan đến điểm đầu và điểm
  cuối, hoặc các đại lượng đại diện tương ứng, chứ không phụ thuộc vào đường
  đi đã dùng để tới điểm đầu, điểm cuối ấy.
  \item Có rất nhiều bài toán con chồng lặp; nếu không, ưu thế của quy hoạch
  động khó thể hiện.
\end{enumerate}

Sau đây ta lấy bài ``Nhận dạng ký tự'' của IOI 1997 để phân tích cách hình
thành suy nghĩ quy hoạch động trong tình huống thông thường.

\paragraph{Bài toán nhận dạng ký tự của IOI 1997.}
Đề bài như sau. Tệp \texttt{FONT.DAT} mô tả dạng chữ của 27 ký hiệu:
khoảng trắng và các chữ cái \texttt{A} đến \texttt{Z}. Ma trận điểm ảnh của
mỗi ký hiệu được biểu diễn bằng 20 dòng, mỗi dòng gồm 20 ký tự \texttt{0} hoặc
\texttt{1}. Một tệp nhập khác mô tả ảnh điểm của một chuỗi ký tự gồm \(N\)
dòng, nhưng các ký tự có thể bị ``hỏng'': một số \texttt{0} biến thành
\texttt{1}, hoặc một số \texttt{1} biến thành \texttt{0}. Mỗi dòng vẫn luôn
có 20 ký tự \texttt{0} hoặc \texttt{1}, nhưng số dòng ứng với mỗi ký tự có thể
rơi vào một trong các trường hợp:
\begin{itemize}
  \item vẫn là 20 dòng, tức không mất dòng nào và không dòng nào bị sao chép;
  \item là 19 dòng, tức có một dòng bị mất;
  \item là 21 dòng, tức có một dòng bị sao chép; hai dòng sao chép có thể bị
  hỏng khác nhau.
\end{itemize}
Yêu cầu là, dưới một giả thiết phân đoạn dòng nào đó, xuất ra kết quả nhận
dạng, tức chuỗi ký tự ứng với phương án làm cho số lần đảo \texttt{0} và
\texttt{1} là ít nhất.

Sau khi sơ bộ xác định bài toán này có thể giải bằng quy hoạch động, ta có thể
dùng phương pháp, hoặc góc nhìn, toán học để mô tả bài toán. Với một bài toán
tối ưu thông thường, vốn cũng là loại bài chính của quy hoạch động, luôn có
một tiêu chuẩn tối ưu hóa. Vì quy hoạch động phải được thực hiện bằng truy
hồi, ta cần phân tích các đại lượng cần thiết để xác định trạng thái. Chẳng
hạn, trong bài nhận dạng ký tự, xét theo quy mô bài toán thì ta đang tìm một
cách phân đoạn và đối ứng cho \(N\) dòng. Vì thế, xét vài dòng đầu là một đại
lượng xác định trạng thái rất tự nhiên. Tiêu chuẩn tối ưu đã được đề bài cho:
dưới một giả thiết nào đó, bao gồm cách phân đoạn và cách nhận dạng, số lần
đảo \texttt{0} và \texttt{1} là ít nhất.

Nếu xem tiêu chuẩn đo lường này là một hàm, thì nó chính là một hàm tối ưu,
hay hàm mục tiêu. Giá trị của hàm tối ưu phụ thuộc vào các biến độc lập, tức
các đại lượng xác định trạng thái. Số biến độc lập phải tùy bài toán; ở đây
chỉ có một biến là số dòng, mang nghĩa ``xét vài dòng đầu''. Điểm then chốt là
xác định trạng thái một cách hữu hiệu. Với trạng thái ấy, các đại lượng đã
chọn phải đủ để xác định giá trị hàm tối ưu; về lý thuyết, khi các đại lượng
ấy đã cố định, hàm tối ưu phải có một giá trị xác định. Nếu không, đại lượng
đã chọn là chưa đủ, hoặc cần dùng đại lượng khác để mô tả. Ngay cả khi đã đủ
xác định, nếu việc lập quan hệ truy hồi gặp khó khăn thì cũng phải điều chỉnh
các biến xác định giá trị hàm tối ưu theo khó khăn ấy. Ngược lại, nếu đặt quá
nhiều đại lượng để xác định hàm tối ưu, hiệu suất quy hoạch động sẽ giảm mạnh:
hoặc phải giải nhiều bài toán con không cần thiết, hoặc biến các bài toán con
chồng lặp thành những bài toán không chồng lặp dưới hệ biến ấy, làm hiệu suất
giảm nghiêm trọng, thậm chí mất hẳn ưu thế của quy hoạch động. Trong ví dụ
này, với \(L\) dòng đầu, hàm tối ưu rõ ràng có giá trị xác định.

Một tư tưởng quan trọng của truy hồi trong quy hoạch động là phân giải bài
toán phức tạp thành các bài toán con. Vì thế, việc xác định không gian bài
toán con và lập quan hệ truy hồi là rất quan trọng. Các biến xác định giá trị
hàm tối ưu thường gợi ý cho không gian bài toán con. Thông thường, bằng cách
suy ngược từ bước gần bài toán nhất, ta có thể thu được một bài toán con có
quy mô nhỏ hơn; lặp lại cách xét ấy, ta thường cảm nhận được không gian bài
toán con của bài toán. Trong quá trình đó, phân tích ngược cũng giúp rút ra
quan hệ truy hồi tương đối dễ dàng. Cần nhấn mạnh rằng đây chỉ là tổng kết
quá trình suy nghĩ khi giải một số bài; về nguyên tắc, các bài khác nhau vẫn
cần được đối xử khác nhau.

Trong bài nhận dạng ký tự, khi xét giá trị hàm tối ưu cho \(n\) dòng, ta chú
ý rằng cuối cùng chắc chắn phải phân đoạn và nhận dạng theo từng ký tự. Do đó,
ký tự cuối cùng chỉ có thể gồm 19, 20, hoặc 21 dòng. Xét lần lượt ba cách
phân đoạn này. Với đoạn vừa cắt ra ứng với một ký tự, trong mỗi phương án cắt
đương nhiên phải chọn ký tự khớp nhất: nếu không dùng ký tự có số lần đảo ít
nhất làm kết quả khớp mà vẫn dẫn tới tối ưu toàn cục, thì chỉ cần thay bước
này bằng ký tự có số lần đảo ít nhất sẽ thu được kết quả tốt hơn ``tối ưu''
đã giả định, mâu thuẫn. Sau khi bỏ đi một ký tự, số dòng giảm xuống; những
dòng còn lại rõ ràng cũng phải được khớp theo cách tối ưu, điều này có thể
chứng minh phản chứng tương tự. Như vậy ta thu được một bài toán con giống
bài toán ban đầu, cùng biến xác định và cùng tiêu chuẩn tối ưu, nhưng quy mô
nhỏ hơn. Đồng thời, quan hệ truy hồi giữa bài toán con và bài toán ban đầu
cũng được thiết lập:
\[
\begin{aligned}
  f[i]=\min\{&
  \mathrm{Compare19}[i-19+1]+f[i-19],\\
  &\mathrm{Compare20}[i-20+1]+f[i-20],\\
  &\mathrm{Compare21}[i-21+1]+f[i-21]\}.
\end{aligned}
\]
Ở đây \(f[i]\) biểu thị giá trị hàm tối ưu khi khớp \(i\) dòng đầu;
\(\mathrm{Compare19}[i]\), \(\mathrm{Compare20}[i]\), và
\(\mathrm{Compare21}[i]\) lần lượt biểu thị số ký tự \texttt{0} và \texttt{1}
bị đảo ít nhất khi 19, 20, hoặc 21 dòng bắt đầu từ dòng \(i\) được khớp với
ký tự gần nhất theo ba cách khớp tương ứng. Điều kiện đầu là \(f[0]=0\); với
số dòng không thể khớp, chỉ cần dùng một số rất lớn đặc biệt để biểu diễn.
Tất nhiên, trọng tâm của bài này không chỉ nằm ở suy nghĩ quy hoạch động cơ
bản. Ví dụ này chỉ dùng để nói về một hướng suy nghĩ cơ bản cho nhiều bài quy
hoạch động; bài còn có phần mô hình hóa toán học như dùng nhị phân để biểu
diễn chuỗi \texttt{0}, \texttt{1}. Mã nguồn xem Chương trình~1 trong phụ lục.

Đôi khi, dù các đại lượng xác định trạng thái thu được theo hướng trên đã đủ
để làm hàm tối ưu có giá trị xác định, việc lập quan hệ truy hồi vẫn gặp khó
khăn. Khi ấy, đưa thêm biến mới hoặc điều chỉnh biến đã có cũng là một cách
khắc phục. Chẳng hạn, xét bài ``Trò chơi xếp khối'' của NOI 1997.

\paragraph{Bài toán xếp khối.}
Khối là một hình hộp chữ nhật. Biết độ dài ba cạnh \(a,b,c\) của \(N\) khối có
đánh số, cần chọn một số khối trong \(N\) khối để xếp thành \(M\) cột
\((1\le M\le N\le 100)\), sao cho tổng chiều cao lớn nhất và thỏa:
\begin{itemize}
  \item mọi khối trong cột thứ \(K\) có số hiệu lớn hơn mọi khối trong cột thứ
  \(K+1\);
  \item với hai khối kề nhau bất kỳ, mặt trên của khối dưới phải chứa được mặt
  dưới của khối trên, và số hiệu khối dưới phải nhỏ hơn số hiệu khối trên.
\end{itemize}

Vì đề yêu cầu cột có số thứ tự nhỏ hơn chứa các khối có số hiệu lớn hơn, điều
này không thật tự nhiên. Không làm thay đổi kết quả, ta đổi thành cột có số
thứ tự nhỏ hơn chứa các khối có số hiệu nhỏ hơn; điều này hiển nhiên không
ảnh hưởng đến tổng chiều cao cuối cùng. Khi đó, mỗi cách xếp hợp lệ có thể
xem là: chọn một số khối theo thứ tự tăng dần của số hiệu, đặt từng cột theo
thứ tự tăng dần của số cột, và trong mỗi cột thì các khối lần lượt được đặt
lên trên khối trước để tạo thành một phương án.

Dùng phương pháp phân tích chung ở trên, ta dễ thấy rằng có thể xét \(i\)
khối đầu được xếp thành \(j\) cột đầu. Hai đại lượng \(i,j\) rõ ràng xác định
được giá trị hàm tối ưu, nhưng quan hệ truy hồi lại không dễ lập trực tiếp.
Phân tích thêm sẽ thấy vấn đề chính nằm ở việc khối thứ \(i\) rốt cuộc có thể
đặt lên cột cuối cùng trong phương án tối ưu của bài toán con, tức trạng thái
xác định bởi \((i-1,j)\), hay không. Nếu xét thêm hai biến là ràng buộc tối
thiểu về dài và rộng của mặt trên khối cuối cùng trong dãy, việc lập truy hồi
không khó. Tuy nhiên, rõ ràng trong quá trình truy hồi có rất nhiều kết quả
không được dùng, làm phình to không gian bài toán con một cách nhân tạo, gây
rắc rối cho lưu trữ và làm hiệu suất thấp.

Thực ra, để lập truy hồi, ta chỉ cần biết mặt trên của cột cuối cùng trong lời
giải bài toán con có chứa được khối hiện tại hay không. Những mặt ràng buộc có
thể xuất hiện trong đề nhiều nhất chỉ có \(3\cdot100+1=301\) loại, gồm cả
trường hợp không ràng buộc. Như vậy, sau khi thêm một đại lượng là ``yêu cầu
mặt trên của cột cuối cùng chứa được loại mặt thứ \(k\)'', quan hệ truy hồi
chỉ cần chia thành ba trường hợp: khối hiện tại mở một cột mới, khối hiện tại
đặt lên cột trước nếu có thể, hoặc không dùng khối hiện tại. Mã nguồn xem
Chương trình~2.

Ngoài ra, có những bài toán vẫn khó lập truy hồi nếu chỉ điều chỉnh quan hệ
truy hồi theo cách trên. Khi ấy, thêm đại lượng hoặc hàm khác để lập công thức
truy hồi cũng là một hướng suy nghĩ, tương tự việc ``tăng cường mệnh đề''
trong chứng minh bằng quy nạp toán học. Một đặc trưng quan trọng khi dùng quy
hoạch động là truy hồi; truy hồi là quá trình dùng kết quả cũ để thu được kết
quả mới. Nếu về lý thuyết có thể chứng minh rằng một bài toán khó truy hồi
trực tiếp có thể được giải cùng lúc với một quan hệ truy hồi mới đưa vào, thì
nhìn qua có vẻ làm bài toán phức tạp hơn, nhưng thực tế ở mỗi bước truy hồi,
ta vừa tăng thêm bài toán cần giải, vừa tăng thêm điều kiện, tức các giá trị
đã giải trước đó, nên truy hồi lại dễ tiến hành hơn.

Ví dụ, xét bài ``Đa giác'' của IOI 1998. Bài toán như sau. Có một đa
giác \(N\) cạnh, trên đỉnh đặt số nguyên, trên cạnh đặt dấu \(+\) hoặc
\(*\). Cần tìm một thứ tự hợp nhất bằng phép toán, qua \(N-1\) phép tính, để
kết quả cuối cùng là lớn nhất.

Nếu chỉ xét \(\mathrm{MAX}[I,L]\), là giá trị lớn nhất có thể thu được khi
bắt đầu từ \(I\) và thực hiện \(L\) phép toán, thì khó truy hồi. Dựa vào kiến
thức toán học, ta đưa thêm \(\mathrm{MIN}[I,L]\), là giá trị nhỏ nhất có thể
thu được khi bắt đầu từ \(I\) và thực hiện \(L\) phép toán. Khi truy hồi, ta
có thể dùng hiệu quả các giá trị với \(I,L\) nhỏ hơn để thu được kết quả lớn
hơn, và thực chất giải đồng thời hai bài toán giá trị nhỏ nhất và lớn nhất.

Quan hệ truy hồi như sau. Xét \(I\) từ \(1\) đến \(N\), \(L\) từ \(1\) đến
\(N-1\). Xét \(t\), vị trí phép toán cuối cùng, từ \(0\) đến \(L-1\). Nếu phép
toán cuối cùng là \(+\), thì:
\[
\begin{aligned}
  \min(i,L) &=
  \min\{\min(i,t)+\min((i+t+1-1)\bmod N+1,L-t-1)\},\\
  \max(i,L) &=
  \max\{\max(i,t)+\max((i+t+1-1)\bmod N+1,L-t-1)\}.
\end{aligned}
\]
Nếu phép toán cuối cùng là \(*\), thì:
\[
\begin{aligned}
  \min(i,L)=\min\{&
  \min(i,t)\min((i+t+1-1)\bmod N+1,L-t-1),\\
  &\min(i,t)\max((i+t+1-1)\bmod N+1,L-t-1),\\
  &\max(i,t)\min((i+t+1-1)\bmod N+1,L-t-1),\\
  &\max(i,t)\max((i+t+1-1)\bmod N+1,L-t-1)\},\\
  \max(i,L)=\max\{&
  \min(i,t)\min((i+t+1-1)\bmod N+1,L-t-1),\\
  &\min(i,t)\max((i+t+1-1)\bmod N+1,L-t-1),\\
  &\max(i,t)\min((i+t+1-1)\bmod N+1,L-t-1),\\
  &\max(i,t)\max((i+t+1-1)\bmod N+1,L-t-1)\}.
\end{aligned}
\]
Mã nguồn xem Chương trình~3.

Ngoài ra, quy hoạch động được thực hiện bằng truy hồi, nên bài toán và bài
toán con càng giống nhau, càng có quy luật thì càng dễ thao tác. Vì thế, với
một số bài mà các giai đoạn và quy luật tự thân không rõ ràng, ta có thể dùng
tiền xử lý để làm chúng đều đặn hơn, dễ cài đặt hơn. Ví dụ, xét bài ``Khớp
biểu thức chính quy'' trong ACM Asia Regional 1997, khu vực Thượng Hải.

\paragraph{Bài toán khớp biểu thức chính quy.}
Biểu thức chính quy là biểu thức có ký tự đại diện. Các ký hiệu mở rộng được
đề định nghĩa gồm:
\begin{itemize}
  \item \texttt{.}: biểu thị ký tự bất kỳ;
  \item \texttt{[c1-c2]}: biểu thị một ký tự bất kỳ nằm giữa \texttt{c1} và
  \texttt{c2};
  \item \texttt{[\^{ }c1-c2]}: biểu thị một ký tự bất kỳ không nằm giữa
  \texttt{c1} và \texttt{c2};
  \item \texttt{*}: ký tự đứng trước nó có thể xuất hiện 0 hoặc nhiều lần;
  \item \texttt{+}: ký tự đứng trước nó có thể xuất hiện một hoặc nhiều lần;
  \item \texttt{\textbackslash}: ký tự đứng sau nó được xem như một ký tự
  thường.
\end{itemize}
Với một chuỗi nhập, cần tìm chuỗi con khớp với biểu thức chính quy và nằm trái
nhất; nếu điều kiện trái nhất như nhau thì chọn chuỗi dài nhất.

Nếu không tiền xử lý, đôi khi một ký hiệu mở rộng có thể ứng với nhiều ký tự,
đôi khi nhiều ký hiệu mở rộng lại chỉ ứng với một ký tự, gây nhiều rắc rối cho
xử lý có hệ thống. Vì vậy cần chuẩn hóa biểu thức chính quy, sao cho hoặc một
nút chỉ ứng với một ký tự, hoặc dùng một đánh dấu đặc biệt để chỉ nó có thể
lặp nhiều lần. Định nghĩa kiểu bản ghi:
\begin{verbatim}
NodeType = Record
  StartChar: Char;  {ky tu bat dau}
  EndChar: Char;    {ky tu ket thuc}
  Belong: Boolean;  {co thuoc khoang hay khong}
  Times: Boolean;   {False: phai xuat hien mot lan;
                     True: co the xuat hien nhieu lan hoac khong xuat hien}
End;
\end{verbatim}

Sau khi tiền xử lý dữ liệu nhập, việc lập quan hệ truy hồi không còn khó.
Dùng \(\mathrm{Pro}[i,j]\) biểu thị việc \(i\) nút biểu thức chính quy đầu có
khớp một chuỗi con kết thúc tại ký tự thứ \(j\) hay không, với giá trị
\(\mathrm{True}/\mathrm{False}\). Với trường hợp \(\mathrm{True}\), đồng thời
ghi rõ vị trí bắt đầu sớm nhất dưới điều kiện này. Nếu nút thứ \(i\) của biểu
thức chính quy chỉ xuất hiện một lần, thì nếu nó không khớp ký tự thứ \(j\),
giá trị là \(\mathrm{False}\); nếu khớp, giá trị giống
\(\mathrm{Pro}[i-1,j-1]\). Khởi đầu đặt \(\mathrm{Pro}[0,x]=\mathrm{True}\).
Nếu nút ấy có thể lặp nhiều lần, nó có thể được giải thích là 0 hoặc nhiều ký
tự. Với điều kiện bản thân nút đó khớp 0 hoặc nhiều ký tự ở các vị trí tương
ứng, ta xét lần lượt các khả năng này; chỉ cần trong đó có \(\mathrm{True}\)
thì \(\mathrm{Pro}[i,j]\) là \(\mathrm{True}\), đồng thời ghi lại vị trí bắt
đầu sớm nhất trong các trường hợp đạt \(\mathrm{True}\). Truy hồi như vậy đến
khi \(i\) đạt số nút. Mã nguồn xem Chương trình~4.

\section{Các vấn đề khi cài đặt quy hoạch động}

Sau khi đã có ý tưởng cơ bản để giải một bài bằng quy hoạch động, nhìn chung
việc cài đặt thuật toán tương đối dễ suy nghĩ, nhưng đôi khi vẫn gặp một số
vấn đề khiến thuật toán khó thực hiện. Về tổng thể, các thuật toán thiết kế
theo tư tưởng quy hoạch động đều truy hồi theo công thức truy hồi đã rút ra.
Đối với máy tính, nếu thiết kế hợp lý, loại truy hồi này thường có hiệu suất
khá cao, nên khả năng quá thời gian không lớn. Ngược lại, quy hoạch động cần
rất nhiều không gian để lưu các kết quả trung gian. Nhờ vậy, mọi bài toán chứa
cùng một bài toán con có thể dùng chung lời giải của bài toán con ấy, thể hiện
ưu thế của quy hoạch động; nhưng ưu thế này phải trả giá bằng không gian. Để
truy cập hiệu quả các kết quả đã có, dữ liệu cũng không dễ nén khi lưu trữ,
nên mâu thuẫn về không gian khá nổi bật.

Mặt khác, hiệu quả thời gian cao của quy hoạch động thường phải được thể hiện
trên dữ liệu kiểm thử lớn, nhất là khi so với tìm kiếm. Vì vậy, với bài toán
quy mô lớn, cách giải quyết vấn đề tràn bộ nhớ mà về cơ bản không ảnh hưởng
đến tốc độ chạy là một vấn đề thường gặp khi dùng quy hoạch động.

Theo tôi, có thể thử suy nghĩ theo các hướng sau.

Một hướng là cố gắng dùng ít không gian nhất. Chẳng hạn, xét từ cấu trúc dữ
liệu của nút, chỉ lưu những nội dung thật sự cần thiết; hoặc tính toán chặt
chẽ phạm vi lưu trữ dữ liệu, như lưu theo bit, lưu nén, v.v. Tất nhiên, việc
này phải phân tích theo từng bài. Ngoài ra, khi cài đặt quy hoạch động, một
cách ta thường dùng là dùng một mảng có số phần tử bằng số nút để lưu quyết
định ở mỗi bước. Điều này rất thuận tiện để suy ngược ra một phương án đạt lời
giải tối ưu, và tốc độ xử lý cũng được cải thiện phần nào. Nhưng khi bộ nhớ
căng thẳng, ta cần nắm mâu thuẫn chính của bài toán: bỏ mảng lưu quyết định
ấy, đổi thành khi suy ngược từng cấp từ lời giải tối ưu, ta tính lại một lần
nữa để chọn một trạng thái ở giai đoạn trước có thể đạt giá trị này, cho đến
khi suy ra kết quả. Cách làm này tốn thêm một chút công viết chương trình so
với cách trước, hiệu quả chạy cũng có thể giảm đôi chút, nhưng thường rất
nhỏ, đổi lại tiết kiệm được nhiều không gian. Vì vậy tư tưởng này có ý nghĩa
trong một số bài toán.

Tuy nhiên, đôi khi dù dùng cách ấy vẫn bị tràn không gian. Lúc đó cần phân
tích xem các dữ liệu được giữ lại có thật sự phải đồng thời nằm trong bộ nhớ
hay không. Trong rất nhiều bài, khi truy hồi quy hoạch động xử lý các phần
phía sau, dữ liệu ở những phần khá xa phía trước thực tế không còn được dùng.
Với loại bài này, ta có thể ghi đè dữ liệu lên những vị trí chắc chắn không
còn dùng nữa, từ đó tái sử dụng không gian. Nếu dùng hữu hiệu cách này, ngay
cả bài toán quy mô khá lớn cũng không đến mức tràn bộ nhớ. Để tìm ra phương
án tối ưu, việc giữ lại quyết định ở từng bước vẫn cần thiết, và điều này
cũng cần không gian.

Nói chung, cách này có thể được cài đặt theo hai hướng. Một là kết quả truy
hồi chỉ dùng hai mảng \(\mathrm{Data1}\) và \(\mathrm{Data2}\). Mỗi lần lấy
\(\mathrm{Data1}\) làm giai đoạn trước, suy ra mảng \(\mathrm{Data2}\), rồi
sao chép \(\mathrm{Data2}\) đè lên \(\mathrm{Data1}\), lặp lại như vậy để thu
được kết quả cuối cùng. Cách này có hạn chế: với bài toán mà truy hồi liên
quan đến vài giai đoạn trước, nó khá phiền phức; hơn nữa mỗi cấp truy hồi lại
phải sao chép nhiều nội dung, tác động càng lớn nếu bài liên quan đến nhiều
giai đoạn trước.

Cách thứ hai là, với bài toán có thể liên quan đến \(N\) giai đoạn trước, lập
mảng \(\mathrm{Data}[0..N]\), trong đó mỗi phần tử có cùng nội dung như
\(\mathrm{Data1}/\mathrm{Data2}\) ở cách trên. Khi đó, nếu không dùng cách
tiết kiệm bộ nhớ thì truy cập chỉ số \(K\); còn khi dùng cách này, chỉ cần
đổi thành truy cập \(K\bmod (N+1)\). So với cách không xử lý như vậy, mã cần
sửa rất ít, tốc độ gần như không bị ảnh hưởng vì phép \(\bmod\) trên máy tính
rất nhanh, và việc giữ lại các số lượng giai đoạn khác nhau cũng dễ thực
hiện. Cách này áp dụng được cho khá nhiều bài, chẳng hạn bài ``Miễn phí bánh
nướng'' của NOI 1998.

\paragraph{Bài toán miễn phí bánh nướng.}
Đề bài như sau. Có một sân khấu rộng \(W\) ô, với \(W\) lẻ và
\(1\le W\le 99\), chiều cao \(H\) ô với \(1\le H\le 100\). Ở thời điểm 0,
người chơi đứng ở chính giữa sân khấu. Mỗi đơn vị thời gian, người chơi có
thể đi sang trái 2 ô, sang trái 1 ô, đứng yên, sang phải 1 ô, hoặc sang phải
2 ô. Mỗi chiếc bánh cho biết thời điểm và vị trí bắt đầu rơi, tốc độ rơi là
số ô đi xuống trong 1 giây, và điểm số. Chỉ khi ở cuối một giây nào đó, bánh
và người chơi ở cùng một ô thì bánh mới được xem là bắt được. Cần tìm một
phương án di chuyển để tổng điểm lớn nhất. Chú ý rằng các bánh đã được sắp
theo thời điểm bắt đầu rơi.

Nhìn vào bài toán, nghĩ đến quy hoạch động không quá khó. Nhưng đề quy định
thời điểm bắt đầu rơi từ 0 đến 1000, và nếu xét đến thời điểm rơi cuối cùng
thì có thể tới khoảng 1100; chiều rộng có thể tới 99. Nếu lấy cặp thời
gian--vị trí làm yếu tố xác định trạng thái để truy hồi thì tốc độ không chậm,
nhưng nếu dùng một mảng cho kết quả sau khi tiền xử lý dữ liệu ban đầu, tức
mảng mô tả tại thời điểm nào, vị trí nào có thể nhận bao nhiêu điểm, một mảng
cho truy hồi quy hoạch động, và một mảng ghi quyết định ở mỗi bước, thì do đề
không chỉ rõ độ lớn điểm số, nếu hai mảng đầu dùng kiểu \texttt{LongInt}, mảng
cuối dùng \texttt{ShortInt}, bộ nhớ cần khoảng \(1100\cdot99\cdot9\) byte,
tức khoảng 957 KB. Rõ ràng điều này không thể lưu nổi.

Nhưng chú ý rằng khi truy hồi, một khi giá trị lớn nhất ứng với một cặp
\((\text{thời gian},\text{vị trí})\) đã xác định, dữ liệu gốc ở vị trí ấy
không còn dùng nữa. Vì vậy hai mảng có thể nhập làm một, chỉ cần
\(1100\cdot99\cdot5\) byte, khoảng 532 KB. Với quy mô của đề, như vậy vừa đủ
để giải quyết. Tất nhiên, nếu suy nghĩ thêm, trong bài này truy hồi thực ra
chỉ liên quan đến thời điểm trước đó, còn dữ liệu bánh chỉ dùng giá trị tại vị
trí hiện tại. Vì thời điểm bắt đầu rơi đã được sắp xếp, khi đọc tới một bánh
có thời điểm bắt đầu muộn hơn thời điểm đang xử lý, không cần đọc ngay. Để
tránh đọc đĩa lặp lại và thiếu quy luật, đồng thời tránh chi phí bộ nhớ quá
lớn, chỉ cần ghi lại tình hình trong khoảng 100 đơn vị thời gian sau thời điểm
hiện tại. Dùng phương pháp tái sử dụng bộ nhớ vòng đã nói ở trên, chỉ cần
\(101\cdot99\cdot4+99\cdot2\cdot2=40392\) byte, chưa tới 40 KB. Với mỗi thời
điểm, chỉ cần 99 số \texttt{ShortInt} để lưu quyết định; dù nâng quy mô bài
toán lên 3000 hoặc 4000 đơn vị thời gian cũng có thể giải thuận lợi. Mã nguồn
xem Chương trình~5.

Khi những phương pháp trên vẫn không giải quyết được vấn đề bộ nhớ, cũng có
thể dùng cấp phát động để khiến phần lớn điểm kiểm thử được giải hiệu quả.
Ngoài ra, dùng bộ nhớ động còn có một lợi ích: khi cần trao đổi trong quá
trình dùng lại bộ nhớ, ta có thể chỉ hoán đổi con trỏ mà không sao chép dữ
liệu. Cách này cũng rất hữu hiệu trong thi đấu.

\section{Kết luận}

Quy hoạch động là một tư tưởng thiết kế chương trình quan trọng. Vì nó không
có hình thức thuật toán cố định nên có tính linh hoạt lớn, nhưng bản chất của
nó lại có mức độ tương đồng rất cao. Vì thế, khi học và sử dụng tư tưởng này,
điểm then chốt là vận dụng linh hoạt trên cơ sở hiểu và nắm chắc bản chất.
Bài viết tuy đã nói đến một số phương pháp tư duy, nhưng đó chỉ là thảo luận
về một số vấn đề khá phổ biến; điều quan trọng nhất vẫn là phân tích cụ thể
từng bài toán cụ thể và lập mô hình toán học phù hợp.

\section*{Tài liệu tham khảo}

\begin{enumerate}
  \item Wu Wenhu, Wang Jiande, \textit{Phân tích thuật toán thực dụng và thiết
  kế chương trình}, Nhà xuất bản Công nghiệp Điện tử, ISBN
  7-5053-4402-1/TP.2036.
  \item Wu Wenhu, Wang Jiande, \textit{Hướng dẫn thi Olympic Tin học quốc tế
  và quốc gia cho thanh thiếu niên: thuật toán và thiết kế chương trình trong
  tổ hợp}, Nhà xuất bản Đại học Thanh Hoa, ISBN 7-302-02203-8/TP.1060.
  \item Đề NOI 1997, NOI 1998, IOI 1997, IOI 1998, và đề ACM Asia Regional
  1997.
\end{enumerate}

\appendix
\section{Chương trình}

\subsection{Chương trình 1: IOI 1997, Nhận dạng ký tự}

Phương trình quy hoạch động cơ bản của bài ``Nhận dạng ký tự'' đã được giải
thích trong phần chính. Ở đây bổ sung điểm then chốt để tăng tốc: khi so sánh
có lệch dòng, cần nâng cao hiệu suất. Trường hợp thiếu một dòng hoặc thừa một
dòng tuy có thể xuất hiện lệch, nhưng mỗi dòng chỉ có thể được so sánh với
một trong hai dòng lân cận. Vì vậy, trước hết ghi lại các phép so sánh có thể
xảy ra, rồi cộng dồn số đảo khi khớp không lệch và khi khớp lệch từ hai đầu
đến một vị trí. Khi xét 20 trường hợp, chỉ cần một vòng lặp ngoài, không tính
vòng lặp trong thủ tục so sánh một dòng, nên hiệu suất tăng đáng kể.

\begin{verbatim}
program Character_Recognition;
const
  cc: array [1..27] of char =
    (' ','A','B','C','D','E','F','G','H','I','J','K','L',
     'M','N','O','P','Q','R','S','T','U','V','W','X','Y','Z');
var
  f, f1, f2: text;
  font: array [1..540] of longint;  {mot longint bieu dien 20 bit}
  dd: array [1..1200] of longint;
  str: string;
  i, j, k: integer;
  t: word;
  ff: integer;
  bin: array [1..20] of longint;
  pro: array [0..1200] of word;
  sta: array [0..1200] of byte;
  bf: array [0..1200] of word;
  pf: array [1..21,0..1] of word;
  n: integer;

procedure getnum(var l: longint);
var
  i: integer;
begin
  l := 0;
  for i := 1 to 20 do
    if str[i] = '1' then inc(l, bin[i]);
end;

function compare0(a, b: longint): byte;
var
  k: byte;
  i: integer;
begin
  a := a xor b;
  k := 0;
  for i := 1 to 20 do
    if a and bin[i] <> 0 then inc(k);
  compare0 := k
end;

function compare20(k: integer; var ff: integer): word;
var
  i, j, t, s: word;
  best: word;
begin
  ff := 0;
  best := maxint;
  for t := 1 to 27 do
  begin
    j := 0;
    for i := 1 to 20 do
    begin
      s := compare0(font[(t-1)*20+i], dd[i+k-1]);
      inc(j, s);
    end;
    if j < best then begin best := j; ff := t end;
  end;
  compare20 := best;
end;

function compare19(k: integer; var ff: integer): word;
var
  i, j, t: word;
  best: word;
  l1, l2: array [0..20] of word;
  bb: integer;
begin
  ff := 0;
  best := maxint;
  for t := 1 to 27 do
  begin
    fillchar(l1, sizeof(l1), 0);
    fillchar(l2, sizeof(l2), 0);
    for i := 1 to 19 do
      for j := 0 to 1 do
        pf[i,j] := compare0(font[(t-1)*20+i+j], dd[k+i-1]);
    l1[1] := pf[1,0];
    for i := 2 to 19 do l1[i] := l1[i-1] + pf[i,0];
    l2[19] := pf[19,1];
    for i := 18 downto 1 do l2[i] := l2[i+1] + pf[i,1];
    bb := maxint;
    for i := 1 to 20 do
      if l1[i-1] + l2[i] < bb then bb := l1[i-1] + l2[i];
    if bb < best then begin best := bb; ff := t end;
  end;
  compare19 := best
end;

function compare21(k: integer; var ff: integer): word;
var
  i, j, t: word;
  best: word;
  l1, l2: array [0..22] of word;
  bb: integer;
begin
  ff := 0;
  best := maxint;
  for t := 1 to 27 do
  begin
    fillchar(l1, sizeof(l1), 0);
    fillchar(l2, sizeof(l2), 0);
    fillchar(pf, sizeof(pf), 0);
    for i := 1 to 21 do
      for j := 0 to 1 do
        if not ((i = 21) and (j = 0)) and
           not ((i = 1) and (j = 1)) then
          pf[i,j] := compare0(font[(t-1)*20+i-j], dd[k+i-1]);
    l1[1] := pf[1,0];
    for i := 2 to 20 do l1[i] := l1[i-1] + pf[i,0];
    l2[22] := 0;
    for i := 21 downto 2 do l2[i] := l2[i+1] + pf[i,1];
    bb := maxint;
    for i := 1 to 20 do
      if l1[i-1] + l2[i+1] + pf[i,0] < bb then
        bb := l1[i-1] + l2[i+1] + pf[i,0];
    if bb < best then begin best := bb; ff := t end;
  end;
  compare21 := best
end;

begin
  assign(f, 'Font.Dat'); reset(f);
  assign(f1, 'Image.dat'); reset(f1);
  assign(f2, 'Image.out'); rewrite(f2);

  readln(f, k);
  bin[1] := 1;
  for i := 2 to 20 do bin[i] := bin[i-1] * 2;

  for i := 1 to k do
  begin
    readln(f, str);
    getnum(font[i]);
  end;

  read(f1, n);
  for i := 1 to n do
  begin
    readln(f1, str);
    getnum(dd[i]);
  end;

  fillchar(pro, sizeof(pro), 255);
  fillchar(sta, sizeof(sta), 0);
  fillchar(bf, sizeof(bf), 255);
  pro[0] := 0; sta[0] := 0; bf[0] := 0;

  for i := 19 to n do
  begin
    if (i >= 19) and (pro[i-19] <> 65535) then
    begin
      t := compare19(i-19+1, ff);
      if t + pro[i-19] < pro[i] then
      begin
        pro[i] := t + pro[i-19];
        sta[i] := ff;
        bf[i] := i - 19;
      end;
    end;
    if (i >= 20) and (pro[i-20] <> 65535) then
    begin
      t := compare20(i-20+1, ff);
      if t + pro[i-20] < pro[i] then
      begin
        pro[i] := t + pro[i-20];
        sta[i] := ff;
        bf[i] := i - 20;
      end;
    end;
    if (i >= 21) and (pro[i-21] <> 65535) then
    begin
      t := compare21(i-21+1, ff);
      if t + pro[i-21] < pro[i] then
      begin
        pro[i] := t + pro[i-21];
        sta[i] := ff;
        bf[i] := i - 21;
      end;
    end;
  end;

  str := '';
  i := n;
  if pro[n] = 65535 then
  begin
    writeln(f2, 'No answer.');
    close(f1); close(f2); halt
  end;
  while i <> 0 do
  begin
    str := cc[sta[i]] + str;
    i := bf[i];
  end;
  writeln(f2, str);
  close(f1); close(f2)
end.
\end{verbatim}

\subsection{Chương trình 2: NOI 1997, Trò chơi xếp khối}

Chương trình dùng truy hồi từng cấp để tránh tràn bộ nhớ.

\begin{verbatim}
program Toy_Bricks_Game;
type
  jmtype = array [0..100] of ^node;
  node = array [0..300] of longint;
var
  f1, f2: text;
  size: array [0..300,1..2] of word;
  num: integer;
  n, m: integer;
  i, j, k, t: integer;
  p, q, r: jmtype;
  aa: array [1..100,1..3] of word;
  ss: array [1..100,1..3] of word;

procedure sort(i: integer);
var
  j, k, t: integer;
begin
  for j := 1 to 2 do
    for k := j+1 to 3 do
      if aa[i,j] > aa[i,k] then
      begin
        t := aa[i,j]; aa[i,j] := aa[i,k]; aa[i,k] := t
      end;
end;

function add(a1, a2: integer): integer;
var
  i: integer;
begin
  for i := 1 to num do
    if (a1 = size[i,1]) and (a2 = size[i,2]) then
    begin add := i; exit end;
  inc(num);
  size[num,1] := a1; size[num,2] := a2;
  add := num;
end;

procedure preprocess;
var
  i: integer;
begin
  num := 0;
  for i := 1 to n do
  begin
    sort(i);
    ss[i,1] := add(aa[i,1], aa[i,2]);
    ss[i,2] := add(aa[i,1], aa[i,3]);
    ss[i,3] := add(aa[i,2], aa[i,3]);
  end;
end;

procedure check(ii, nn, hh: integer);
begin
  if (p[ii-1]^[0] >= 0) and
     (p[ii-1]^[0] + hh >= q[ii]^[j]) then
    q[ii]^[j] := p[ii-1]^[0] + hh;
  if (p[ii-1]^[ss[ii,nn]] >= 0) and
     (q[ii-1]^[ss[ii,nn]] + hh >= q[ii]^[j]) then
    q[ii]^[j] := q[ii-1]^[ss[ii,nn]] + hh;
end;

begin
  assign(f1, 'ToyBrick.in'); reset(f1);
  assign(f2, 'ToyBrick.out'); rewrite(f2);
  readln(f1, n, m);
  for i := 1 to n do readln(f1, aa[i,1], aa[i,2], aa[i,3]);

  size[0,1] := 0; size[0,2] := 0;
  preprocess;

  for i := 0 to n do
  begin
    new(p[i]); new(q[i]);
    fillchar(q[i]^, sizeof(q[i]^), 0);
  end;

  for t := 1 to m do
  begin
    r := q; q := p; p := r;
    for i := 0 to n do fillchar(q[i]^, sizeof(q[i]^), 255);
    for i := t to n do
      for j := 0 to num do
      begin
        q[i]^[j] := q[i-1]^[j];
        if (aa[i,1] >= size[j,1]) and (aa[i,2] >= size[j,2]) then
          check(i, 1, aa[i,3]);
        if (aa[i,1] >= size[j,1]) and (aa[i,3] >= size[j,2]) then
          check(i, 2, aa[i,2]);
        if (aa[i,2] >= size[j,1]) and (aa[i,3] >= size[j,2]) then
          check(i, 3, aa[i,1]);
      end;
  end;
  writeln(f2, q[n]^[0]);
  close(f1); close(f2)
end.
\end{verbatim}

\subsection{Chương trình 3: IOI 1998, Đa giác}

\begin{verbatim}
program Polygon;
var
  f1, f2: text;
  n: integer;
  data: array [1..50] of integer;
  sign: array [1..50] of char;
  i, j, k, l: integer;
  t, s, p: integer;
  ans: set of 1..50;
  best: integer;
  min, max: array [1..50,0..50] of integer;
  first: boolean;

procedure init;
var
  i: integer;
  ch: char;
begin
  readln(f1, n);
  for i := 1 to n do
  begin
    repeat read(f1, ch) until ch <> ' ';
    sign[i] := ch;
    read(f1, data[i]);
  end;
end;

begin
  assign(f1, 'Polygon.in'); reset(f1);
  assign(f2, 'Polygon.out'); rewrite(f2);
  init;
  best := -maxint - 1;
  ans := [];
  fillchar(max, sizeof(max), 0);
  fillchar(min, sizeof(min), 0);
  for j := 1 to n do
  begin
    max[j,0] := data[j];
    min[j,0] := data[j];
  end;

  for l := 1 to n-1 do
    for k := 1 to n do
    begin
      max[k,l] := -maxint - 1;
      min[k,l] := maxint;
      for t := 0 to l-1 do
        case sign[(k+t+1-1) mod n + 1] of
          't':
          begin
            s := max[k,t] + max[(k+t+1-1) mod n + 1,l-t-1];
            if s > max[k,l] then max[k,l] := s;
            s := min[k,t] + min[(k+t+1-1) mod n + 1,l-t-1];
            if s < min[k,l] then min[k,l] := s;
          end;
          'x':
          begin
            for p := 1 to 4 do
            begin
              case p of
                1: s := max[k,t] * max[(k+t+1-1) mod n + 1,l-t-1];
                2: s := max[k,t] * min[(k+t+1-1) mod n + 1,l-t-1];
                3: s := min[k,t] * max[(k+t+1-1) mod n + 1,l-t-1];
                4: s := min[k,t] * min[(k+t+1-1) mod n + 1,l-t-1];
              end;
              if s > max[k,l] then max[k,l] := s;
              if s < min[k,l] then min[k,l] := s;
            end;
          end;
        end;
    end;

  for i := 1 to n do
    if max[i,n-1] > best then begin best := max[i,n-1]; ans := [i] end
    else if max[i,n-1] = best then include(ans, i);

  writeln(f2, best);
  first := true;
  for i := 1 to n do
    if i in ans then
    begin
      if first then first := false else write(f2, ' ');
      write(f2, i);
    end;
  writeln(f2);
  close(f1); close(f2)
end.
\end{verbatim}

\subsection{Chương trình 4: ACM 1997 khu vực châu Á, Thượng Hải: khớp biểu thức chính quy}

\begin{verbatim}
program Expression_Match;
type
  datatype = record
    st, ed: char;
    md: 0..1;  {0: mot lan; 1: lap 0 hoac nhieu lan}
    mt: 0..1;  {0: khop khi khong thuoc; 1: khop khi thuoc}
  end;
var
  f1, f2: text;
  s1, s2: string;
  len: integer;
  dd: array [1..80] of datatype;
  pro: array [0..80,0..80] of boolean;
  fr: array [0..80,0..80] of byte;
  i, j, k: integer;
  ok: boolean;
  bt, bj: integer;

procedure preprocess;
var
  i, j: integer;
begin
  i := 0; j := 0;
  while i < length(s1) do
  begin
    inc(i);
    case s1[i] of
      '.':
      begin
        inc(j);
        dd[j].md := 0; dd[j].mt := 1;
        dd[j].st := #0; dd[j].ed := #255;
      end;
      '*': dd[j].md := 1;
      '+':
      begin
        inc(j);
        dd[j] := dd[j-1];
        dd[j].md := 1;
      end;
      '\':
      begin
        inc(i); inc(j);
        dd[j].md := 0; dd[j].mt := 1;
        dd[j].st := s1[i]; dd[j].ed := s1[i];
      end;
      '[':
      begin
        inc(i); inc(j);
        dd[j].md := 0; dd[j].mt := 1;
        if s1[i] = '^' then begin inc(i); dd[j].mt := 0 end;
        if s1[i] = '\' then inc(i);
        dd[j].st := s1[i];
        inc(i, 2);
        if s1[i] = '\' then inc(i);
        dd[j].ed := s1[i];
        inc(i);
      end
      else
      begin
        inc(j);
        dd[j].st := s1[i]; dd[j].ed := s1[i];
        dd[j].mt := 1; dd[j].md := 0;
      end;
    end;
  end;
  len := j
end;

begin
  assign(f1, 'Match.in'); reset(f1);
  assign(f2, 'Match.out'); rewrite(f2);
  while true do
  begin
    readln(f1, s1);
    if s1 = 'end' then break;
    readln(f1, s2);
    preprocess;
    ok := false;
    fillchar(pro, sizeof(pro), false);
    fillchar(pro[0], sizeof(pro[0]), true);
    fillchar(fr, sizeof(fr), 0);
    bt := maxint; bj := 0;
    for i := 0 to length(s2) do fr[0,i] := i + 1;
    for i := 1 to len do
      for j := 1 to length(s2) do
      begin
        if dd[i].md = 0 then
        begin
          if (dd[i].mt = 0) xor (s2[j] in [dd[i].st..dd[i].ed]) then
          begin
            pro[i,j] := pro[i-1,j-1];
            if pro[i,j] then fr[i,j] := fr[i-1,j-1];
          end
          else pro[i,j] := false;
        end
        else
        begin
          for k := j downto 0 do
          begin
            if pro[i-1,k] then
            begin
              pro[i,j] := true;
              if (fr[i,j] = 0) or (fr[i,j] > fr[i-1,k]) then
                fr[i,j] := fr[i-1,k];
            end;
            if not ((dd[i].mt = 0) xor (s2[k] in [dd[i].st..dd[i].ed])) then
              break;
          end;
        end;
        if (i = len) and pro[i,j] and (fr[i,j] <= bt) then
        begin
          ok := true;
          bt := fr[i,j];
          bj := j - bt + 1;
        end;
      end;

    if ok then
      if bj <> 0 then
        writeln(f2, copy(s2,1,bt-1), '(', copy(s2,bt,bj), ')',
                copy(s2,bt+bj,length(s2)-bt-bj+1))
      else writeln(f2, '()', s2)
    else writeln(f2, s2);
  end;
  close(f1); close(f2)
end.
\end{verbatim}

\subsection{Chương trình 5: NOI 1998, Miễn phí bánh nướng}

Phương trình quy hoạch động là
\[
  \mathrm{ans}[t,j]=
  \max_{k\in\{-2,-1,0,1,2\}}\mathrm{ans}[t-1,j+k],
\]
trong đó vị trí tương ứng phải đạt được.

\begin{verbatim}
program Pizza_For_Free;
type
  link = ^linetype;
  linetype = array [1..99] of shortint;
var
  f1, f2: text;
  w, h: integer;
  dd: array [0..100] of array [1..99] of longint;
  pro: array [0..3100] of link;
  dt: array [0..1,1..99] of longint;
  i, j, k, t: integer;
  bt, bj: integer;
  best: longint;
  maxt: integer;
  ans: array [1..3100] of shortint;
  _t, _j, _v: integer;
  _fz: longint;
  _saved: boolean;

procedure try_reading;
var
  t0, j, v: integer;
  fz: longint;
begin
  fillchar(dd[(t+100) mod 101], sizeof(dd[(t+100) mod 101]), 0);
  if _saved then
  begin
    if _t > t then exit;
    _saved := false;
    if (h-1) mod _v = 0 then
    begin
      inc(dd[(_t+(h-1) div _v) mod 101,_j], _fz);
      if (_t+(h-1) div _v) > maxt then maxt := _t+(h-1) div _v;
    end;
  end;
  if eof(f1) then exit;

  while true do
  begin
    if eof(f1) then exit;
    readln(f1, t0, j, v, fz);
    if t0 > t then
    begin
      _v := v; _t := t0; _j := j; _fz := fz;
      _saved := true;
      exit
    end;
    if (h-1) mod v <> 0 then continue;
    inc(dd[(t0+(h-1) div v) mod 101,j], fz);
    if t0+(h-1) div v > maxt then maxt := t0+(h-1) div v;
  end;
end;

begin
  assign(f1, 'INPUT.TXT'); reset(f1);
  assign(f2, 'OUTPUT.TXT'); rewrite(f2);
  readln(f1, w, h);

  for i := 0 to 3100 do
  begin
    new(pro[i]);
    fillchar(pro[i]^, sizeof(pro[i]^), 0);
  end;
  for i := 0 to 1 do
    for j := 1 to 99 do
      dt[i,j] := -maxlongint - 1;

  dt[0,(1+w) div 2] := 0;
  fillchar(dd, sizeof(dd), 0);
  t := 0; maxt := 0; _saved := false;
  try_reading;
  best := dd[0,(w+1) div 2];
  bt := 0; bj := (w+1) div 2;

  while true do
  begin
    inc(t);
    try_reading;
    if eof(f1) and (t > maxt) and not _saved then break;
    for j := 1 to w do
    begin
      dt[t mod 2,j] := -maxlongint - 1;
      for k := -2 to 2 do
        if (j+k > 0) and (j+k <= w) and
           (dt[1-t mod 2,j+k] > -maxlongint - 1) and
           (dt[1-t mod 2,j+k] + dd[t mod 101,j] > dt[t mod 2,j]) then
        begin
          dt[t mod 2,j] := dt[1-t mod 2,j+k] + dd[t mod 101,j];
          pro[t]^[j] := k;
        end;
      if dt[t mod 2,j] > best then
      begin
        best := dt[t mod 2,j];
        bt := t; bj := j;
      end;
    end;
  end;

  writeln(f2, best);
  j := bj;
  for t := bt downto 1 do
  begin
    ans[t] := -pro[t]^[j];
    inc(j, pro[t]^[j]);
  end;
  for t := 1 to bt do writeln(f2, ans[t]);
  close(f1); close(f2)
end.
\end{verbatim}

\section{Nội dung chi tiết của các bài toán được dẫn}

\subsection{Nhận dạng ký tự (IOI 1997)}

Bài toán yêu cầu viết một chương trình thực hiện nhận dạng ký tự.

Mỗi ảnh ký tự lý tưởng có 20 dòng, mỗi dòng 20 chữ số; mỗi chữ số là
\texttt{0} hoặc \texttt{1}. Tệp \texttt{FONT.DAT} chứa biểu diễn của 27 ảnh
ký tự lý tưởng theo thứ tự: khoảng trắng, rồi \texttt{a} đến \texttt{z}.

Tệp \texttt{IMAGE.DAT} chứa một hoặc nhiều ảnh ký tự có thể bị hỏng. Một ảnh
ký tự có thể hỏng theo các cách sau:
\begin{itemize}
  \item nhiều nhất một dòng bị sao chép, và dòng sao chép đứng ngay sau dòng
  gốc;
  \item nhiều nhất một dòng bị thiếu;
  \item một số ký tự \texttt{0} bị đổi thành \texttt{1};
  \item một số ký tự \texttt{1} bị đổi thành \texttt{0}.
\end{itemize}
Không có ảnh ký tự nào vừa bị sao chép một dòng vừa bị thiếu một dòng. Trong
dữ liệu đánh giá, không quá 30\% số \texttt{0} và \texttt{1} của bất kỳ ảnh
ký tự nào bị đổi. Khi một dòng bị sao chép, một hoặc cả hai dòng kết quả đều
có thể bị hỏng, và hỏng theo các cách khác nhau.

Nhiệm vụ là nhận dạng chuỗi gồm một hoặc nhiều ký tự trong ảnh cho bởi
\texttt{IMAGE.DAT}, dùng phông chữ cho bởi \texttt{FONT.DAT}. Khi nhận dạng
một ảnh ký tự, hãy chọn ảnh ký tự trong phông chữ cần số lần đổi tổng cộng
giữa \texttt{0} và \texttt{1} ít nhất để trở thành ảnh đã cho, dưới giả thiết
thuận lợi nhất về dòng bị sao chép hoặc bị thiếu. Trong trường hợp dòng bị
sao chép, chỉ tính lỗi trên dòng ít hỏng hơn. Mọi ký tự trong ảnh mẫu và ảnh
đánh giá đều có thể được một chương trình tốt nhận dạng từng ký tự một. Mỗi
bộ dữ liệu đánh giá có nghiệm tốt nhất duy nhất.

Một lời giải đúng phải dùng chính xác toàn bộ dữ liệu trong tệp
\texttt{IMAGE.DAT}.

\paragraph{Dữ liệu vào.}
Cả hai tệp nhập đều bắt đầu bằng một số nguyên \(N\), với
\(19\le N\le 1200\), chỉ số dòng tiếp theo:
\begin{verbatim}
N
(digit1)(digit2)(digit3)...(digit20)
(digit1)(digit2)(digit3)...(digit20)
...
\end{verbatim}
Mỗi dòng dữ liệu rộng 20 chữ số. Không có dấu cách giữa các số 0 và 1. Tệp
\texttt{FONT.DAT} mô tả phông chữ, luôn chứa 541 dòng, và có thể khác nhau
giữa các bộ kiểm thử.

\paragraph{Kết quả.}
Chương trình phải tạo tệp \texttt{IMAGE.OUT}, chứa một chuỗi duy nhất là các
ký tự đã nhận dạng, trên một dòng văn bản ASCII. Kết quả không chứa ký tự phân
cách. Nếu chương trình không nhận dạng được một ký tự nào đó, phải in
\texttt{?} ở vị trí tương ứng. Định dạng kết quả nói trên thay thế yêu cầu
chuẩn trong quy chế về việc dùng dấu cách phân tách trong kết quả.

Điểm số được tính theo tỉ lệ ký tự nhận dạng đúng. Trong mẫu của đề, tệp
\texttt{IMAGE.DAT} mô tả một ký tự \texttt{A} bị hỏng và kết quả mẫu là:
\begin{verbatim}
IMAGE.OUT
A
\end{verbatim}

\subsection{Trò chơi xếp khối (NOI 1997)}

SERCOI gần đây thiết kế một trò chơi xếp khối. Mỗi người chơi có \(N\) khối
hộp chữ nhật được đánh số lần lượt \(1,2,3,4,\ldots\). Ba cạnh khác nhau của
mỗi khối lần lượt gọi là cạnh \(a\), cạnh \(b\), và cạnh \(c\).

Luật chơi như sau:
\begin{enumerate}
  \item Chọn một số khối trong \(N\) khối và chia chúng thành \(M\) cột
  \((1\le M\le N)\), gọi là cột 1, cột 2, cột 3, \(\ldots\). Mỗi cột có ít
  nhất một khối, và mọi khối trong cột thứ \(K\) phải có số hiệu lớn hơn mọi
  khối trong cột thứ \(K+1\).
  \item Với mỗi cột, người chơi xếp các khối thẳng đứng thành một trụ, thỏa
  hai điều kiện: trừ khối trên cùng, mặt trên của mỗi khối tiếp xúc với đúng
  một mặt dưới của khối khác; mặt trên của khối dưới phải chứa được mặt dưới
  của khối trên, tức độ dài hai cặp cạnh của mặt trên khối dưới lần lượt lớn
  hơn hoặc bằng hai cặp cạnh của mặt dưới khối trên. Ngoài ra, với hai khối
  tiếp xúc trên dưới bất kỳ, số hiệu khối dưới phải nhỏ hơn số hiệu khối trên.
\end{enumerate}
Cuối cùng, thắng thua được quyết định bởi tổng chiều cao của \(M\) trụ đã xếp.
Hãy viết chương trình tìm một phương án xếp để tổng chiều cao của \(M\) trụ là
lớn nhất.

\paragraph{Dữ liệu vào và kết quả.}
Tệp nhập là \texttt{INPUT.TXT}. Dòng đầu có hai số nguyên dương \(N\) và
\(M\), \((1\le M\le N\le 100)\), lần lượt là số khối và số trụ cần xếp. Hai
số cách nhau bởi một dấu cách. Tiếp theo là \(N\) dòng theo thứ tự số hiệu
khối từ 1 đến \(N\); mỗi dòng có ba số nguyên từ 1 đến 1000, là độ dài các
cạnh \(a,b,c\) của khối ấy.

Tệp ra là \texttt{OUTPUT.TXT}, chỉ có một dòng là một số nguyên, biểu thị tổng
chiều cao của \(M\) trụ. Mẫu trong bản gốc được trích từ bố cục PDF và các
dấu cách trong dữ liệu bị mất một phần; về ý nghĩa, nó minh họa một tệp nhập
với \(N=4\), \(M=2\), các kích thước khối theo sau, và kết quả là 24.

\subsection{Đa giác (IOI 1998)}

Trò chơi đa giác là trò chơi một người. Ban đầu cho một đa giác có \(N\) đỉnh;
mỗi đỉnh được gán một giá trị nguyên, mỗi cạnh được gán một ký hiệu \(+\) hoặc
\(*\). Tất cả các cạnh được đánh số từ 1 đến \(N\).

Nước đi đầu tiên cho phép xóa một cạnh. Mỗi nước đi tiếp theo gồm:
\begin{enumerate}
  \item chọn một cạnh \(E\), cùng hai đỉnh \(V_1\) và \(V_2\) nối bởi \(E\);
  \item dùng một đỉnh mới thay thế cạnh \(E\) và hai đỉnh \(V_1,V_2\). Đỉnh
  mới được gán giá trị bằng kết quả áp dụng phép toán ghi trên \(E\) cho
  \(V_1\) và \(V_2\).
\end{enumerate}
Khi mọi cạnh đều bị xóa, chỉ còn một đỉnh và trò chơi kết thúc. Điểm của trò
chơi là giá trị của đỉnh còn lại. Phần ví dụ chơi trong bản gốc chỉ ghi là
không in tại đây.

Nhiệm vụ là viết chương trình, với một đa giác bất kỳ, tính điểm cao nhất có
thể đạt được và liệt kê mọi cạnh có thể là cạnh bị xóa ở nước đầu tiên để dẫn
đến điểm cao nhất ấy.

Tệp \texttt{POLYGON.IN} gồm hai dòng. Dòng đầu chứa \(N\). Dòng thứ hai chứa
ký hiệu được gán cho từng cạnh \(1,\ldots,N\), cùng các giá trị đỉnh xen giữa
hai cạnh. Giá trị nguyên đầu tiên ứng với đỉnh nối đồng thời cạnh 1 và cạnh
2; giá trị thứ hai ứng với đỉnh nối cạnh 2 và cạnh 3; cứ tiếp tục như vậy,
giá trị cuối cùng ứng với đỉnh nối cạnh \(N\) và cạnh 1. Ký hiệu và giá trị
cách nhau bằng một dấu cách. Ký hiệu cạnh có hai loại: chữ \texttt{t} ứng với
\(+\), và chữ \texttt{x} ứng với \(*\).

Ví dụ nhập:
\begin{verbatim}
4
t -7 t 4 x 2 x 5
\end{verbatim}
Tệp này ứng với đa giác trong hình của đề; ký tự đầu tiên ở dòng hai là ký
hiệu của cạnh 1.

Trong tệp \texttt{POLYGON.OUT}, dòng đầu in điểm cao nhất có thể đạt được.
Dòng thứ hai liệt kê theo thứ tự tăng dần mọi cạnh mà nếu xóa ở nước đầu tiên
có thể dẫn đến điểm cao nhất, các số cách nhau bởi một dấu cách.

Ví dụ kết quả:
\begin{verbatim}
33
1 2
\end{verbatim}

\subsection{Khớp biểu thức chính quy (ACM 1997 khu vực châu Á, Thượng Hải)}

Biểu thức chính quy là một chuỗi chứa các ký tự thường và một số ký tự meta.
Các ký tự meta gồm:
\begin{itemize}
  \item \texttt{.}: ký tự bất kỳ;
  \item \texttt{[c1-c2]}: ký tự bất kỳ nằm giữa \texttt{c1} và \texttt{c2};
  \item \texttt{[\^{ }c1-c2]}: ký tự bất kỳ không nằm giữa \texttt{c1} và
  \texttt{c2};
  \item \texttt{*}: ký tự trước nó có thể xuất hiện số lần bất kỳ, kể cả
  không xuất hiện;
  \item \texttt{+}: ký tự trước nó có thể xuất hiện số lần bất kỳ nhưng ít
  nhất một lần;
  \item \texttt{\textbackslash}: mọi ký tự theo sau được xem là ký tự thường.
\end{itemize}

Cần viết chương trình tìm chuỗi con nằm trái nhất của một chuỗi cho trước sao
cho chuỗi con ấy khớp với biểu thức chính quy đã cho. Nếu có nhiều chuỗi con
cùng khớp và có cùng vị trí bắt đầu bên trái, chọn chuỗi dài nhất.

\paragraph{Dữ liệu vào.}
Mỗi hai dòng của tệp nhập là một bài toán khớp mẫu. Dòng đầu là biểu thức
chính quy, dòng sau là chuỗi cần khớp. Mỗi dòng dài không quá 80 ký tự. Một
dòng chỉ chứa \texttt{end} kết thúc dữ liệu vào.

\paragraph{Kết quả.}
Với mỗi bài toán, in một dòng. Dòng này chứa chuỗi cần khớp, nhưng chuỗi con
trái nhất khớp với biểu thức chính quy phải được đặt trong ngoặc tròn. Nếu
không có chuỗi con nào khớp, in nguyên chuỗi nhập.

Ví dụ nhập:
\begin{verbatim}
.*
asdf
f.*d.
sefdfsde
[0-9]+
asd345dsf
[^\*-\*]
**asdf**fasd
b[a-z]*r[s-u]*
abcdefghijklmnopqrstuvwxyz
[T-F]
dfkgjf
end
\end{verbatim}

Ví dụ kết quả:
\begin{verbatim}
(asdf)
se(fdfsde)
asd(345)dsf
**(a)sdf**fasd
a(bcdefghijklmnopqrstu)vwxyz
dfkgjf
\end{verbatim}

\subsection{Miễn phí bánh nướng (NOI 1998)}

SERKOI vừa giới thiệu một trò chơi có tên ``Miễn phí bánh nướng''.

Trò chơi diễn ra trên một sân khấu. Sân khấu rộng \(W\) ô, màn trời cao
\(H\) ô, và người chơi chiếm một ô. Ban đầu người chơi đứng ở chính giữa sân
khấu, tay cầm một chiếc khay. Sau khi trò chơi bắt đầu, bánh liên tục xuất
hiện từ các ô ở mép trên màn trời và rơi thẳng xuống. Người chơi di chuyển
trái phải để hứng bánh. Mỗi giây người chơi có thể sang trái hoặc sang phải
một ô hoặc hai ô, cũng có thể đứng yên.

Có nhiều loại bánh. Người chơi đã chấm điểm từng loại bánh theo khẩu vị của
mình. Dưới sự điều khiển từ xa của máy 8-308, tốc độ rơi của các loại bánh
cũng khác nhau, tính bằng ô mỗi giây. Nếu ở cuối một giây nào đó, bánh vừa
đến đúng ô người chơi đang đứng, người chơi thu được bánh ấy.

Hãy viết chương trình giúp người chơi thu bánh sao cho tổng điểm của các bánh
thu được là lớn nhất.

\paragraph{Dữ liệu vào.}
Dòng đầu của tệp nhập gồm hai số nguyên dương cách nhau bởi dấu cách: chiều
rộng \(W\) của sân khấu, là số lẻ từ 1 đến 99, và chiều cao \(H\), từ 1 đến
100. Các dòng tiếp theo, theo thứ tự thời điểm bắt đầu rơi của bánh, cho
thông tin về từng chiếc bánh. Mỗi dòng gồm bốn số nguyên dương: thời điểm bắt
đầu rơi, từ 0 đến 1000 giây; vị trí ngang; tốc độ rơi, từ 1 đến 100; và điểm
số. Thời điểm trò chơi bắt đầu là 0. Các ô theo chiều ngang được đánh số từ
trái sang phải, bắt đầu từ 1. Các mục trên cùng một dòng được cách nhau bởi
một hoặc nhiều dấu cách.

\paragraph{Kết quả.}
Dòng đầu của tệp ra in một số nguyên dương là tổng điểm lớn nhất chương trình
có thể thu được. Mỗi dòng tiếp theo, theo thứ tự thời gian, in quyết định của
người chơi trong từng giây. In \texttt{0} nghĩa là đứng yên; \texttt{1} hoặc
\texttt{2} nghĩa là sang phải một hoặc hai ô; \texttt{-1} hoặc \texttt{-2}
nghĩa là sang trái một hoặc hai ô. Kết quả kéo dài đến khi người chơi thu xong
chiếc bánh cuối cùng mà phương án chọn sẽ thu.

Ví dụ nhập:
\begin{verbatim}
3 3
0 1 2 5
0 2 1 3
1 2 1 3
1 3 1 4
\end{verbatim}

Ví dụ kết quả:
\begin{verbatim}
12
-1
1
1
\end{verbatim}

\end{document}
